\section{The Depth Division of Labor}
\label{sec:motivation}
This section presents the paper's central finding---an \emph{understand-then-generate}
division of labor along a transformer's depth---and derives CoMem's design from
it. We organize it as seven numbered insights. The first three (I1--I3) are
\emph{direct, method-independent measurements} of where understanding lives:
layer-wise probes, causal layer-truncation, and split-depth recomputation
sweeps. They stand on their own as a claim about transformer computation---the
lower layers finish understanding early and the top layers only generate---and
motivate a new kind of memory that stores that \emph{comprehension} state rather
than tokens or KV. Insight I4 records the cost of caching deeper, and I5--I7
establish how the cached state must be read back (retrieval, made iterative for
chains, with full cross-chunk attention). Numbers marked \emph{oracle} or
single-pass recur as controlled ablations in \S\ref{sec:experiments}.

\paragraph{(I1) Semantic content saturates at mid-depth, so a single mid-layer
hidden is a sufficient statistic for ``understanding'' and is what should be
cached.} Linear probes on semantic tasks peak in the middle of the stack and
fall off at the top, whereas next-token ability forms only in the top layers;
the two curves separate cleanly and consistently across Qwen3-8B and Llama-3.
Causally, truncating the network to only its first $\sim$4--8 layers (normalized
depth $0.12$--$0.22$) already recovers $95\%$ of the full model's downstream
semantics, and an RTE mid-layer probe beats the top-layer probe by
$+0.06$ to $+0.10$. This mid-depth ceiling is exactly where the cache stops
helping: a zero-training split-depth sweep on RULER \citep{ruler}
niah\_single at 16k holds recall at $100$ for every $j\le 9$. The same shape
survives at radically larger scale: on the 80-layer Hunyuan Hy3 MoE the split
sweep shows a flat plateau over $j\in[8,32]$ (top-1 agreement $\approx0.90$,
$\mathrm{KL}\approx0.21$), placing the cacheable ceiling near $0.4L$. Because a
mid-depth residual hidden encodes the chunk's semantics but not yet its
continuation \citep{stagesofinference}, we \textbf{split the transformer at
depth $j$ and \textsc{Write} the depth-$j$ residual hidden}
$h_j\in\mathbb{R}^{512\times d}$ ($\sim$$1/(2L)$ of a full KV cache),
recomputing only layers $[j{:}L]$ at read time.

\paragraph{(I2) The upper layers perform query-conditioned generation, so their
outputs cannot be cached query-blind---the cache point must sit \emph{below} the
split.} The upper band is where the query gets injected, and a hidden computed
without the query diverges sharply from its query-aware counterpart. Measuring
this ``top-prepay'' divergence---computing the top $b$ layers query-blind
instead of recomputing them with the query present---the gap grows monotonically
with $b$: $b{=}6\to$ cosine $0.916$ / relative-$L_2$ $4.2$;
$b{=}8\to 0.902 / 5.8$; $b{=}12\to 0.865 / 11.3$. A direction-B negative
control that caches a \emph{top}-layer hidden confirms the cost end-to-end:
BABILong \citep{babilong} qa5 falls from $(12,0){=}61/50$ (caching below the
split) to $(12,6){=}29/20$ and $(6,12){=}9/10$ as the cache creeps into the
generation band. The trade is visible even within a single layer: the top
layer's native RTE verbalizer scores $0.79$ while its linear probe scores only
$0.62$---the top trades linear separability for generation utility. The same
signature reappears on Hy3 as a ``fidelity smile'': resuming from too \emph{deep}
a split ($j\gtrsim 36$) makes the LM tax rise again, because a deep chunk-local
encoding has already committed to a query-blind continuation. We therefore make
\textsc{Read} pack \texttt{[sink; selected $h_j$; query $h_j$]} and
\textbf{recompute the query-conditioned upper layers with the query present},
keeping the cache point strictly below the generation band.

\paragraph{(I3) The cacheable depth is a bounded knee: $j$ is a
RAG$\leftrightarrow$closed-book knob whose sweet spot lands at
$\sim$$0.33$--$0.40\,L$.} The two endpoints are fixed by construction: $j{=}0$
is full-RAG recompute, which equals the full forward to
$\max|\Delta\text{logit}|<10^{-4}$, and $j{=}L$ is closed-book. Between them the
usable region is a narrow knee. The zero-training niah\_single 16k sweep is flat
at $100$ through $j\le 9$, then falls off a cliff to $14$ at $j{=}12$ and to $0$
at $j{=}18$; a self-distilled LoRA (\S\ref{sec:method}) pushes the usable
ceiling from $j\approx9$ to $j\ge12$. Under perfect retrieval the decline is
smooth and confirms the same band: BABILong qa5 oracle recall at
$j{=}0/6/12/18$ is $69/50/39/16$. Across backbones the knee is stable in
\emph{relative} depth---the dense-8B operating point sits near $0.33L$
($j{=}12$ of $L{=}36$) and the
Hy3 knee at $j\approx32/80\,({=}0.40L)$---even though the absolute
semantic depth is backbone-dependent. We therefore \textbf{operate at
$j\approx0.33$--$0.40\,L$ (e.g.\ $j{=}12$ on Qwen3-8B), expose $j$ as an
explicit knob, and self-distill to push the usable split deeper.}

\paragraph{(I4) Caching deeper costs monotonically more LM quality, so $j$ is a
\emph{trade-off}, not a tax-minimum---and the cache point can be made more
compressible by design (severable).} Forcing a low-rank funnel at successive
depths of a from-scratch 1B model (dim $512$) shows the readout LM tax rises
monotonically with depth: bottleneck at layer $1/3/6/9/12$ costs
$+4.2/5.8/6.0/9.5/9.6\%$ perplexity; narrowing the funnel at layer $6$
($d{=}1024/512/256$) costs $+4.5/5.9/8.5\%$; and every arm is cheaper on the
larger 3B model than on 1B. So the shallow end minimizes LM tax while the deep
end maximizes cacheable semantics (I1, I3), and we therefore
\textbf{set $j$ at this compromise (e.g.\ $j{=}12$ on Qwen3-8B), not at a
free-lunch tax minimum}. We frame this with a prediction-not-reconstruction
information-bottleneck view: the cache need only preserve the predictive
information $I(X;Y_{\text{pred}})$, not the reconstruction information
$I(X;X)$, which is why a single mid-layer hidden suffices and why a
token-reconstruction auxiliary was the wrong objective. This
compressibility-by-design story is a \emph{severable} companion direction
developed in follow-up work: CoMem itself requires
\emph{no} special pretraining and runs on a stock backbone, so we carry only the
tax numbers here as evidence that $j$ is a trade-off.

\paragraph{(I5) At long range, facts drown without retrieval, so the read must
retrieve and recompute; the long-range bottleneck is retrieval quality, not
compression fidelity.} Stripping retrieval isolates the point. HCache-style
mid-layer recompute \emph{without} retrieval \citep{hcache} collapses on long
needle tasks---niah\_single $8\text{k}{=}34\to16\text{k}{=}2$, LongEval all-$0$,
BABILong qa1 near-zero ($16/3/0/0$ at 4k/8k/16k/32k)---yet short,
redundant-fact tasks survive without retrieval (qa5 $2\text{k}{=}68$,
$8\text{k}{=}52$). The failure is therefore \emph{retrieval-missing}, not model
incapacity: the needle is present in the cache but drowned. Keeping the
\emph{recent} tokens instead of the \emph{relevant} ones fails the same way:
StreamingLLM \citep{streamingllm} at the same fixed $\sim$6.7k-token budget
scores niah\_single $90/42/16/4$ at $8/16/64/128$k versus CoMem's
$100/100/100/100$ ($25\times$ at 128k; Table~\ref{tab:slm}). Conversely, given
the right chunk, the compressed read-out is lossless---an oracle selector scores
$100\%$ on NIAH at every length---so what remains to be solved is selection, not
fidelity. We therefore \textbf{feed the read pack from a BM25 top-$k$ selector}
and recompute with full cross-chunk attention (I7).

\paragraph{(I6) Multi-hop reference chains need \emph{iterative} retrieval---a
single pass drowns the chain.} A one-shot selector cannot recover an $N$-hop
assignment chain, because the chunk holding the final value is useless without
the chunks that link to it. On RULER variable-tracking a single pass at 16k
scores only bm25 $27.6$ / reader\_attn $60.2$ / oracle $9.2$ (the oracle is
ill-defined here---locating the answer chunk misses the chain), and all selectors
collapse to $23$ or below at 32k. Making retrieval iterative---\texttt{iter\_bm25}
follows the variable chain (retrieve, read the referenced name, re-retrieve)
with no forward pass and no training---lifts variable-tracking to
$8\text{k}{=}95.2$, $16\text{k}{=}93.8$, $32\text{k}{=}96.8$ (hop-4, $k{=}16$)
and, crucially, stops it degrading with length (Table~\ref{tab:itervt}). The
gain is scoped to reference chains: on independent multi-key needles iterative
retrieval slightly hurts ($92$ vs.\ single-pass $97$--$99$) and on single needles
it is a no-op ($100{=}100$). We therefore add \textbf{\texttt{iter\_bm25} as a
selector variant} (\S\ref{sec:method}), applied to variable-tracking and
multi-hop queries but not to independent-needle tasks.

\paragraph{(I7) Cross-chunk full-attention recompute over the pack is
load-bearing for multi-fact disambiguation.} Because the read recomputes the
upper layers, it can either attend across all retrieved chunks or reuse each
chunk's KV block-diagonally; the difference decides multi-fact tasks. With full
cross-chunk attention, RULER niah\_multikey scores $88/92$ at 16k/32k versus
$44/40$ block-diagonal ($\Delta{+}44/{+}52$); on BABILong the pattern repeats
(qa2 $36/24$ vs.\ $16/12$; qa5 $68/65$ vs.\ $49/53$), while single-needle tasks
are unaffected ($100{=}100$, since one chunk already answers). Disambiguating
which of several candidate facts satisfies the query requires the query to
attend jointly over all of them. We therefore \textbf{recompute layers $[j{:}L]$
with full attention over the entire read pack}, not block-diagonal KV-reuse; the
quantitative ablation is Table~\ref{tab:crosschunk}.
